\begin{tikzpicture}[
  every node/.style={font=\tiny},
  state/.style={draw=institucionalBlue, rounded corners=3pt, fill=institucionalGrey!35, align=center, minimum width=1.86cm, minimum height=0.56cm, line width=0.65pt},
  critical/.style={draw=institucionalRed, rounded corners=3pt, fill=institucionalRed!8, align=center, minimum width=1.86cm, minimum height=0.56cm, line width=0.75pt},
  arrow/.style={-{Latex[length=1.35mm]}, draw=institucionalBlue, line width=0.62pt},
  arrowred/.style={-{Latex[length=1.35mm]}, draw=institucionalRed, line width=0.75pt}
]
\foreach \i/\txt in {0/Creada,1/Planeación,2/Lista,3/Ejecución,4/Terminada,5/Informe,6/Acta,7/SES,8/Factura,9/Pago,10/Cerrada}{
  \ifnum\i>6
    \node[critical] (s\i) at (\i*1.76,0) {\txt};
  \else
    \node[state] (s\i) at (\i*1.76,0) {\txt};
  \fi
}
\foreach \i in {0,...,6}{\pgfmathtruncatemacro{\j}{\i+1}\draw[arrow] (s\i)--(s\j);}
\foreach \i in {7,...,9}{\pgfmathtruncatemacro{\j}{\i+1}\draw[arrowred] (s\i)--(s\j);}
\node[below=0.72cm of s5, align=center, font=\scriptsize, text=institucionalBlue] {Cada transición requiere precondiciones, permisos y evidencia documental};
\end{tikzpicture}
